\corrigesChapitre{2}{Atomes, noyaux et ions}
\par\noindent\begin{minipage}{\linewidth}
\begin{corrige}{1}{Constituants}
\begin{questions}
\item Dans le noyau.
\item L’électron.
\end{questions}
\end{corrige}
\end{minipage}
\par\noindent\begin{minipage}{\linewidth}
\begin{corrige}{2}{Neutralité}
\begin{questions}
\item Huit électrons, car il est neutre.
\item $Z=8$.
\end{questions}
\end{corrige}
\end{minipage}
\par\noindent\begin{minipage}{\linewidth}
\begin{corrige}{3}{Un noyau}
\begin{questions}
\item $Z=11$.
\item $A=11+12=23$.
\end{questions}
\par\smallskip\noindent\textbf{Pour comprendre.}\enspace Le numéro atomique compte les protons : $Z=11$. Le nombre de nucléons est $A=11+12=23$. Un atome neutre possède autant d’électrons que de protons.\par
\pcFigureEx{figures/rev-cor-atome-na.pdf}{Inventaire des constituants ; représentation symbolique sans échelle de taille ni orbites électroniques.}
\end{corrige}
\end{minipage}
\par\noindent\begin{minipage}{\linewidth}
\begin{corrige}{4}{L’eau}
\begin{questions}
\item Deux H et un O.
\item Six H et trois O.
\end{questions}
\pcFigureEx{figures/rev-cor-trois-eaux.pdf}{Le coefficient multiplie la molécule entière : trois molécules contiennent six H et trois O.}
\end{corrige}
\end{minipage}
\par\noindent\begin{minipage}{\linewidth}
\begin{corrige}{5}{Le signe}
\begin{questions}
\item Il a gagné un électron.
\item Il a perdu un électron.
\end{questions}
\pcFigureEx{figures/rev-cor-ions.pdf}{Perdre un électron donne une charge positive ; en gagner un donne une charge négative. Le noyau reste inchangé.}
\end{corrige}
\end{minipage}
\par\noindent\begin{minipage}{\linewidth}
\begin{corrige}{6}{Un ion cuivre}
\begin{questions}
\item $29-2=27$ électrons.
\item 29 protons ; le noyau n’a pas changé.
\end{questions}
\par\smallskip\noindent\textbf{Pour comprendre.}\enspace Le noyau de cuivre conserve ses 29 protons. La charge $2+$ signifie deux électrons de moins que l’atome neutre : $29-2=27$. Ne pas retrancher des protons.\par
\pcFigureEx{figures/rev-cor-atome-cu.pdf}{Inventaire des constituants ; représentation symbolique sans échelle de taille ni orbites électroniques.}
\end{corrige}
\end{minipage}
\par\noindent\begin{minipage}{\linewidth}
\begin{corrige}{7}{Un ion chlorure}
\begin{questions}
\item 18 électrons.
\item Non : une entité monoatomique chargée diffère d’une molécule neutre à deux atomes.
\end{questions}
\pcFigureEx{figures/rev-cor-atome-cl.pdf}{Inventaire des constituants ; représentation symbolique sans échelle de taille ni orbites électroniques.}
\end{corrige}
\end{minipage}
\par\noindent\begin{minipage}{\linewidth}
\begin{corrige}{8}{Inventaire complet}
\begin{questions}
\item 14 neutrons et 13 électrons.
\item Charge $3+$ ; $\ce{Al^{3+}}$.
\end{questions}
\par\smallskip\noindent\textbf{Pour comprendre.}\enspace $N=A-Z=27-13=14$ neutrons. L’atome neutre possède 13 électrons ; l’ion $\ce{Al^3+}$ en a $13-3=10$.\par
\pcFigureEx{figures/rev-cor-atome-al.pdf}{Inventaire des constituants ; représentation symbolique sans échelle de taille ni orbites électroniques.}
\end{corrige}
\end{minipage}
\par\noindent\begin{minipage}{\linewidth}
\begin{corrige}{9}{Indices et coefficient}
\begin{questions}
\item Deux C et huit H.
\item Six N et trois O.
\end{questions}
\pcFigureEx{figures/rev-cor-coefficients.pdf}{Compter chaque type d’atome en multipliant coefficient et indice.}
\end{corrige}
\end{minipage}
\par\noindent\begin{minipage}{\linewidth}
\begin{corrige}{10}{Trois écritures}
\begin{questions}
\item O : symbole de l’oxygène ou atome selon le contexte ; dioxygène : molécule à deux O ; ion oxyde : un O chargé $2-$.
\item Oui, ils ont le même numéro atomique pour chaque noyau d’oxygène.
\end{questions}
\pcFigureEx{figures/rev-cor-oxygene.pdf}{Atome, molécule et ion sont trois entités différentes ; la charge ne se lit pas comme un nombre d’atomes.}
\end{corrige}
\end{minipage}
\par\noindent\begin{minipage}{\linewidth}
\begin{corrige}{11}{Retrouver un élément}
\begin{questions}
\item Deux charges positives non compensées : $2+$.
\item $\ce{Mg^{2+}}$.
\end{questions}
\par\smallskip\noindent\textbf{Pour comprendre.}\enspace La charge en unités élémentaires vaut $12-10=+2$ ; avec $Z=12$, il s’agit de $\ce{Mg^2+}$. Le signe vient du déficit d’électrons.\par
\pcFigureEx{figures/rev-cor-atome-mg.pdf}{Inventaire des constituants ; représentation symbolique sans échelle de taille ni orbites électroniques.}
\end{corrige}
\end{minipage}
\par\noindent\begin{minipage}{\linewidth}
\begin{corrige}{12}{Une fausse justification}
\begin{questions}
\item Un gain ou une perte d’électron ne change pas l’élément. Un ion chlorure provient du chlore ; un ion sodium reste du sodium.
\item Le nombre de protons, donc Z ; cela sort de la transformation chimique usuelle étudiée.
\end{questions}
\end{corrige}
\end{minipage}
\par\noindent\begin{minipage}{\linewidth}
\begin{corrige}{13}{Un changement d’échelle}
\begin{questions}
\item $10^{-10}/10^{-15}=10^5$.
\item $10/10^5=10^{-4}\unite m=0{,}1\unite{mm}$.
\end{questions}
\par\smallskip\noindent\textbf{Pour comprendre.}\enspace $10^{-10}/10^{-15}=10^5$. Si l’atome mesure 10 m dans une maquette, le noyau mesurerait $10/10^5=10^{-4}\unite m=0{,}1\unite{mm}$.\par
\pcFigureEx{figures/rev-cor-echelles-atome.pdf}{Deux zooms séparés : le diamètre du noyau est environ cent mille fois plus petit que celui de l’atome.}
\end{corrige}
\end{minipage}
\par\noindent\begin{minipage}{\linewidth}
\begin{corrige}{14}{Une solution électriquement neutre}
\begin{questions}
\item La somme des charges est nulle.
\item Quatre ions chlorure : $2\times2-4=0$.
\end{questions}
\par\smallskip\noindent\textbf{Pour comprendre.}\enspace Chaque ion chlorure porte la charge $-1$. Deux ions calcium portent ensemble $2\times(+2)=+4$ : il faut donc quatre ions chlorure pour obtenir $+4-4=0$.\par
\pcFigureEx{figures/rev-cor-neutralite.pdf}{Additionner les charges de tous les ions ; conserver les signes et les coefficients.}
\end{corrige}
\end{minipage}
\par\noindent\begin{minipage}{\linewidth}
\begin{corrige}{15}{Deux noyaux du même élément}
\begin{questions}
\item Respectivement 6 et 8 neutrons.
\item Leur même nombre de protons les rattache au carbone ; ce sont deux isotopes, terme fourni ici.
\end{questions}
\pcFigureEx{figures/rev-cor-atome-isotopes.pdf}{Inventaire des constituants ; représentation symbolique sans échelle de taille ni orbites électroniques.}
\end{corrige}
\end{minipage}
\par\noindent\begin{minipage}{\linewidth}
\phantomsection\label{cor-defi-2}\pcSolutionDefi{2}{L’entité inconnue}\begin{questions}
\item $Z=26$ et $A=56$.
\item $\ce{Fe^{3+}}$, car trois électrons manquent par rapport à l’atome neutre.
\item Le nombre de protons reste 26 : l’élément est le fer. La charge vient du déséquilibre entre protons et électrons. L’entité décrite est un ion monoatomique, sans assemblage de plusieurs atomes.
\end{questions}
\par\smallskip\noindent\textbf{Pour comprendre.}\enspace $Z=26$, $A=26+30=56$. La charge vaut $26-23=+3$, donc $\ce{Fe^3+}$ ; les neutrons ne contribuent pas à cette charge.\par
\pcFigureEx{figures/rev-cor-atome-fe3.pdf}{Inventaire des constituants ; représentation symbolique sans échelle de taille ni orbites électroniques.}
\end{minipage}
\par\noindent\begin{minipage}{\linewidth}
\phantomsection\label{cor-activite-2}\pcSousTitre{Pistes pour l’investigation}
Fixer Z et construire un noyau ; l’atome neutre a Z électrons. Retirer deux électrons donne un ion 2+, en ajouter un donne un ion 1−. Les protons ne changent pas. La taille des jetons, leurs distances et leurs positions fixes ne décrivent pas le nuage réel. Le tableau périodique sert à vérifier le symbole et le numéro atomique.\par\end{minipage}
